% !TeX root = ../../tesis.tex

\subsection{Tecnología del servicio de datos}
\label{subsec:backend}

El servicio de \texttt{Apimetro} está construido en \textbf{Go}, lenguaje
compilado con el rendimiento necesario para atender consultas concurrentes sobre
el grafo metropolitano. La arquitectura interna —framework web, patrón de diseño
y configuración de entorno Docker— se documenta en el Anexo~
\ref{anx:apimetro-backend}.

\subsubsection{El modelo SaaS: por qué Apimetro no procesa dos veces}
\label{subsubsec:estructuraEndpoints}

La decisión de arquitectura más relevante de \texttt{Apimetro} no es tecnológica
sino metodológica: ¿debe cada herramienta que necesita datos de la red
procesarlos desde cero, o conviene que exista un servicio centralizado que los
prepare una sola vez?

El modelo adoptado es el del
\textbf{\termino{Software as a Service} (\textsc{saas})}: \texttt{Apimetro}
ingesta los archivos \gtfs{} oficiales, aplica el proceso de normalización y
almacena el resultado. A partir de ese punto, cualquier cliente —
\textit{VFTModel}, un visor web, un script de análisis— consulta la \api{} REST
y recibe geometrías con atributos operativos ya integrados
\citep{cabrera2026VFTModel}. Sin este principio, cada ejecución de
\textit{VFTModel} requeriría releer los mismos archivos de la \textsc{adip},
multiplicando el tiempo de cómputo y abriendo margen a inconsistencias entre
corridas.

La \api{} expone tres tipos de consulta que mapean directamente a las tres
entidades del modelo relacional:

\begin{itemize}
    \item \textbf{Estaciones:} puntos georreferenciados filtrables por
    sistema, alcaldía o jerarquía nodal.
    \item \textbf{Líneas:} trazos con métricas operativas integradas
    (velocidad, frecuencia, capacidad del vehículo).
    \item \textbf{Polígonos:} límites administrativos para el filtrado por
    demarcación territorial.
\end{itemize}

El formato de estas respuestas —y la razón de preferir GeoJSON sobre GeoPackage—
se desarrolla en la sección~\ref{subsec:estructura-geojson}.
